\documentclass[12pt,a4paper]{article}
\usepackage[ngerman]{babel}
\usepackage[utf8]{inputenc}
\usepackage[T1]{fontenc}
\usepackage{geometry}
\usepackage{amsmath}
\usepackage{graphicx}
\usepackage{hyperref}
\usepackage{enumitem}
\usepackage{csquotes}

\geometry{margin=2.5cm}
\hypersetup{hidelinks}

\title{Künstliche Intelligenz als Herausforderung für Polizei und Justiz in Deutschland\\
	\large Eine Analyse rechtlicher, ethischer und praktischer Problemfelder}
\author{}
\date{\today}

\begin{document}
	
	\maketitle
	
	\begin{abstract}
		Die zunehmende Nutzung Künstlicher Intelligenz (KI) durch Polizei und Justiz in Deutschland eröffnet ein breites Spektrum neuer Herausforderungen. Die vorliegende Analyse systematisiert diese Problemfelder, von rechtlichen Grundlagendefiziten über das Risiko algorithmischer Diskriminierung bis hin zu konkreten Gefahren von Fehlurteilen durch KI-Halluzinationen und verzerrte Entscheidungsgrundlagen. Ein besonderer Fokus liegt auf dem neu entstehenden strukturellen Problem: dem strategischen Wettbewerbsvorteil privater Kanzleien gegenüber öffentlichen Institutionen, der die Fairness des Rechtsstaats gefährdet.
	\end{abstract}
	
	\section{Einleitung}
	
	Die Integration Künstlicher Intelligenz in die polizeiliche und justizielle Arbeit schreitet in Deutschland rasch voran. Dieser technologische Wandel bringt jedoch tiefgreifende Herausforderungen mit sich. Auf dem 3. IPA-Fachforum im Mai 2025 in Stuttgart brachte der IPA-Präsident die zentrale Fragestellung auf den Punkt: \enquote{Wie viel KI braucht unser Rechtsstaat – und wie viel hält er aus?} \cite{IPA2025}
	
	\section{Problemfelder beim polizeilichen Einsatz von KI}
	
	\subsection{Rechtliche Defizite und mangelnde Transparenz}
	
	Viele KI-Systeme im polizeilichen Kontext werden ohne ausreichende Rechtsgrundlage, Transparenz und Risikoabschätzung entwickelt, getestet oder genutzt \cite{Peteranderl2025}. Diese grundlegende Kritik wird durch die jüngste Debatte um die \enquote{Lex Palantir} bestätigt, bei der Bürgerrechtsgruppen einen \enquote{unverhältnismäßigen Eingriff} durch automatisierte Datenanalysen beklagen \cite{FR2026}.
	
	Das Bundesverfassungsgericht hat in seiner \enquote{Hessendata}-Entscheidung vom Februar 2023 verfassungsrechtliche Leitplanken für automatisierte Datenanalysen formuliert und die neuartige Eingriffsintensität dieser Technologie betont \cite{VerfBlog2026}. Die Konferenz der unabhängigen Datenschutzaufsichtsbehörden (DSK) forderte im September 2025 eine verfassungskonforme Ausgestaltung dieser Verfahren und betonte die Wahrung digitaler Souveränität \cite{DSK2025}.
	
	\subsection{Risiko algorithmischer Diskriminierung}
	
	KI-Systeme können strukturelle Vorurteile aus Trainingsdaten übernehmen und systematisch gegen bestimmte Gruppen diskriminieren. Die Forschungsstelle der Antidiskriminierungsstelle des Bundes warnt, dass der Einsatz von KI in der Polizei \enquote{sehr riskant} sei, da Algorithmen mitentscheiden könnten, welche Stadtteile verstärkt überwacht werden \cite{Antidisk2025}. \cite{taz2026} weisen darauf hin, dass Algorithmen Vorurteile wie Rassismus reproduzieren und Diskriminierung sogar verstärken können.
	
	Zivilgesellschaftliche Organisationen kritisieren insbesondere, dass automatisierte Auswertungen sich nicht auf Tatverdächtige beschränken, sondern auch Unbeteiligte erfassen und marginalisierte Gruppen dadurch besonders belasten können \cite{Netzpolitik2025}.
	
	\subsection{Ausufernde Überwachung und Grundrechtseingriffe}
	
	Der Einsatz KI-gestützter Systeme wie Gesichtserkennung in Echtzeit ermöglicht eine flächendeckende Überwachung des öffentlichen Raums. Im Frankfurter Bahnhofsviertel wurde im Juli 2025 ein entsprechendes Pilotprojekt gestartet \cite{TOnline2025}. Der Deutsche Anwaltverein warnt, dass biometrische Gesichtserkennung einen \enquote{schweren Eingriff in das Recht auf informationelle Selbstbestimmung} darstelle \cite{DAV2025}.
	
	Der Deutschlandfunk berichtet von Plänen des Bundesinnenministeriums, die Befugnisse für Gesichtserkennung massiv auszubauen, einschließlich des Abgleichs mit Bildern aus dem Internet \cite{DLF2025}. Ein Gutachten von Algorithm Watch zeigt zudem, dass polizeiliche Foto-Fahndung im Internet nur mit illegalen Datenbanken funktionieren kann \cite{taz2025}.
	
	Die \enquote{Lex Palantir} wird von Amnesty International als \enquote{rote Linie} kritisiert, die die Infrastruktur zur flächendeckenden Überwachung der Bevölkerung schaffe \cite{Amnesty2026}.
	
	\subsection{Das Blackbox-Problem und fehlende Wirksamkeitsbelege}
	
	Die oft undurchschaubaren Entscheidungswege von KI-Systemen können fatale Folgen haben. Der Deutsche Anwaltverein bemängelt, dass die Funktionsweise der Algorithmen völlig intransparent sei \cite{DAV2025}. \cite{Peteranderl2025} betonen, dass Entwicklung und Einsatz häufig ohne unabhängige wissenschaftliche Begleitung sowie Kontrollmechanismen erfolgen – obwohl die Systeme erhebliche Risiken bergen.
	
	Besonders problematisch ist, dass es für prädiktive Polizeiarbeit (\enquote{Predictive Policing}) in Deutschland bislang keinen wissenschaftlichen Nachweis für deren Wirksamkeit gibt \cite{Peteranderl2025}. \cite{CILIP2025} bestätigen diese Einschätzung: Es existiere kein wissenschaftlicher Beleg für den Einfluss von Predictive Policing auf eine sinkende Kriminalitätsrate.
	
	\subsection{Abhängigkeit von US-Tech-Konzernen}
	
	Die Nutzung der Analysesoftware \enquote{Gotham} des US-Unternehmens Palantir führt zu einer problematischen technologischen Abhängigkeit. In Baden-Württemberg löste der Alleingang von Innenminister Strobl beim Vertragsabschluss mit Palantir im März 2025 einen Koalitionskonflikt aus \cite{FR2026}. Kritiker befürchten Hintertüren für US-Geheimdienste \cite{Focus2025}, obwohl Palantir dies zurückweist \cite{Spiegel2025}. In Hessen wird die Software bereits seit 2017 eingesetzt, musste aber nach einem Urteil des Bundesverfassungsgerichts angepasst werden \cite{FR2026}.
	
	Einige Bundesländer wie Niedersachsen und Bremen setzen bewusst nicht auf Palantir \cite{Goslarsche2025}. In Baden-Württemberg strebt die Koalition eine europäische Alternative bis 2030 an \cite{FR2026}.
	
	\section{Problemfelder in der Justiz}
	
	\subsection{KI als Hilfsmittel, nicht als Richterersatz}
	
	Der Thüringer Richterbund zieht eine klare Grenze beim KI-Einsatz in der Justiz. Der Landesvorsitzende Holger Pröbstel erklärte gegenüber der Deutschen Presse-Agentur: \enquote{Was aus meiner Sicht gar nicht geht, ist, dass Sie sich von ChatGPT die Urteile schreiben lassen} \cite{SZ2025}. Die Findung eines Urteils müsse immer Aufgabe eines Richters sein: \enquote{Sonst können Sie die Justiz gleich ganz abschaffen} \cite{SZ2025}. KI dürfe für Gerichte und Staatsanwaltschaften immer nur ein Hilfsmittel sein, niemals den Richter ersetzen \cite{SZ2025}.
	
	\subsection{Gefährdung der richterlichen Unabhängigkeit}
	
	Eine tiefe inhaltliche Einbindung von KI in die Entscheidungsfindung kann die richterliche Unabhängigkeit untergraben. Aus der bisherigen Rechtsprechung ist abzuleiten, dass verpflichtende Vorgaben zur Nutzung von KI-Tools im richterlichen Kernbereich unzulässig sind \cite{Anwaltsblatt2025}.
	
	\subsection{Fehleranfälligkeit und mangelnde Standards}
	
	Skeptisch zeigt sich Pröbstel insbesondere gegenüber Überlegungen, mittels KI den Strafrahmen zu bestimmen. Die entscheidende Frage sei: \enquote{Welche Urteile stehen in dieser Datenbank und wie bewerten das die Algorithmen} \cite{SZ2025}. Da die Justiz auf diese Algorithmen keinen Zugriff habe, könne sie so nicht seriös arbeiten \cite{SZ2025}.
	
	Bisher fehlen zudem einheitliche Standards für den Umgang mit KI-generierten Beweismitteln. Die automatisierte Gesichtserkennung in der Strafverfolgung wird in Deutschland bereits seit rund zwei Jahrzehnten praktiziert – allerdings ohne taugliche Rechtsgrundlage \cite{Nomos2025}.
	
	\subsection{Fehlerhafte Verfahrenshandlungen durch KI-Halluzinationen}
	
	Eine der konkretesten Gefahren für Fehlurteile geht vom unkritischen Einsatz generativer KI in der anwaltlichen und richterlichen Arbeit aus. Große Sprachmodelle neigen dazu, plausible, aber frei erfundene Informationen zu produzieren. Eine im Juni 2025 veröffentlichte Datenbank des französischen Anwalts Damien Charlotin weist allein für das Jahr 2025 bereits über 80 vor Gericht aufgedeckte Fälle von KI-Halluzinationen nach \cite{Charlotin2025}.
	
	\subsubsection{Praxisbeispiel Amtsgericht Köln (2025)}
	Ein Rechtsanwalt reichte einen mit Hilfe von KI verfassten Schriftsatz ein, der nicht existierende Gerichtsurteile zitierte. Das Gericht rügte den Anwalt öffentlich und erkannte den Schriftsatz nicht als vollwertig an. Der Fall zeigt, wie ein Mangel an der anwaltlichen Sorgfaltspflicht das gesamte Verfahren beeinträchtigen kann \cite{LTO2025a}.
	
	Weltweit wurden inzwischen mehr als 1.200 gerichtliche Sanktionen gegen Rechtsanwälte wegen fehlerhafter KI-Nutzung verhängt, davon rund 800 allein in den USA \cite{NPR2025}. Einige Gerichte verlangen mittlerweile die Offenlegung der KI-Nutzung in Schriftsätzen.
	
	\subsubsection{Risiken für die Urteilsfindung}
	Wenn Richter Entscheidungen auf KI-Empfehlungen stützen, ohne die Algorithmen zu verstehen, entsteht eine \enquote{Blackbox-Justiz}. Studien aus dem Jahr 2025 zeigen, dass bis zu 15 Prozent der automatisierten Fallanalysen durch KI-Systeme fehlerhaft sind \cite{LTO2025b}. Eine groß angelegte Untersuchung der Europäischen Rundfunkunion stellte zudem eine Fehlerquote von bis zu 40 Prozent bei Antworten gängiger Chatbots fest \cite{heise2025}.
	
	\subsection{Manipulierte oder unzuverlässige Beweise als Grundlage von Fehlurteilen}
	
	Die zunehmende Qualität von Deepfakes ermöglicht die Manipulation audiovisueller Beweismittel. Im November 2025 reichten Kläger in einem US-amerikanischen Zivilprozess ein Video einer angeblichen Zeugin als Beweis ein, das sich später als Deepfake entpuppte. Die Richterin wies die Klage ab, nachdem sie technische Artefakte identifiziert hatte \cite{LTO2025c}. Auch wenn es sich um einen US-Fall handelt, verdeutlicht er das erhebliche Gefahrenpotenzial für deutsche Gerichte.
	
	Eine weitere Beweisrisikoquelle liegt in der unzulässigen Nutzung von KI durch gerichtlich bestellte Sachverständige. Das Landgericht Darmstadt stellte am 10. November 2025 klar, dass ein Sachverständigengutachten, das umfangreich mit KI erstellt wurde, als Beweismittel unzulässig ist. Das Gericht monierte, dass der Gutachter persönliche Untersuchungen unterlassen und sich statt auf die Akten auf KI-generierte Fallzusammenfassungen gestützt hatte, die nicht mit den Tatsachen übereinstimmten. Der Gutachter erhielt keinerlei Vergütung \cite{LG Darmstadt2025}. Solche mangelhaften Gutachten können im Zweifel zu Fehlurteilen beitragen.
	
	\subsection{Verzerrte Entscheidungsgrundlagen durch algorithmische Verzerrungen}
	
	Fehlurteile können auch durch inhärente Verzerrungen in den Trainingsdaten von KI-Systemen entstehen. Wenn ein Algorithmus aus historischen, möglicherweise diskriminierenden Gerichtsurteilen lernt, wird er diese Fehler wiederholen und potenziell verstärken. Ein Forschungsprojekt an der Universität Konstanz untersucht derartige \enquote{algorithmic biases} und ihre Auswirkungen auf die Strafjustiz \cite{Uni Konstanz2025}. Diese systematischen Verzerrungen sind besonders tückisch, da sie nicht auf einen einzelnen menschlichen Fehler zurückgehen, sondern flächendeckend in die Strafzumessung einfließen können.
	
	\subsection{Wettbewerbsvorteil privater Kanzleien durch KI}
	
	Über die unmittelbaren Gefahren von Fehlurteilen hinaus zeichnet sich ein strukturelles Problem ab: Private Kanzleien nutzen KI-Technologie strategisch und oft schneller als öffentliche Institutionen, was zu einem massiven \enquote{Ungleichgewicht der Waffen} (asymmetry of arms) führen kann.
	
	\subsubsection{Strategische Frühphase und unternehmerischer Druck}
	Kanzleien erkennen den Wettbewerbsvorteil durch KI und investieren gezielt. Eine im Oktober 2025 veröffentlichte Benchmark-Studie des Unternehmens Vals AI, die auf einem Datensatz von über 200 juristischen Fragestellungen aus acht US-Kanzleien basierte, zeigte, dass vier getestete KI-Produkte bei den meisten Aufgaben besser abschnitten als die menschliche Vergleichsgruppe. Die durchschnittliche Genauigkeit der KI-Produkte lag bei 80\%, während die der Jurist:innen 71\% erreichte. Besonders bei der Überprüfbarkeit von Quellen (\enquote{Belegbarkeit}) lagen spezialisierte Legal-KI-Produkte mit 76\% deutlich vor menschlichen Jurist:innen (68\%). Noch deutlicher war der Unterschied bei der Mandantenfreundlichkeit und Verständlichkeit (\enquote{Angemessenheit}) – hier erreichten die Jurist:innen nur 60\%, während die KI-Produkte auf 70\% kamen \cite{ValAI2025}.
	
	Prof. Dr. Frauke Rostalski, Direktorin des Instituts für Strafrecht und Strafprozessrecht an der Universität zu Köln, sieht darin ein grundlegendes Problem für die Fairness des Rechtsstaats. Sie betont, dass KI-Systeme für Privatpersonen vor allem eine unterstützende Rolle im Bereich der rechtlichen Selbsthilfe spielen können. Das daraus entstehende Machtgefälle könnte Bürger gegenüber einem übermächtigen Staat jedoch weiter benachteiligen, wenn der Staat hier nicht Schritt hält \cite{Rostalski2024}.
	
	\subsubsection{Strukturelle Ungleichgewichte in der Prozessführung}
	Dieser Effizienzvorteil wird in Massenverfahren besonders deutlich. Der Einsatz von KI bei der Dokumentenanalyse (Technology Assisted Review) ermöglicht es Kanzleien, riesige Datenmengen – tausende E-Mails, Verträge, Gutachten – extrem schnell zu durchsuchen und prozessrelevante Informationen zu extrahieren. Es ist zudem zunehmend üblich, KI-Tools wie \enquote{Harvey} als \enquote{Gedankenpartner} für die Entwicklung einer Prozessstrategie zu nutzen \cite{Hengeler2025}. Eine KI kann tausende relevante Urteile analysieren, Argumentationsmuster erkennen und dem gegnerischen Anwalt, der sich manuell durch die Akten kämpfen muss, einen entscheidenden Schritt voraus sein.
	
	\subsubsection{Konkrete Machtasymmetrie: Fallbeispiel}
	Ein konkretes Fallbeispiel verdeutlicht diese Asymmetrie: Eine Partei verfügt über Tausende von Schadensberichten zu einzelnen beschädigten Gegenständen. Das Unternehmen nutzte eine KI-Plattform, um Beweise vorzubereiten, indem es die Dokumente hochlud und die KI Fragen dazu stellen ließ. Die Gegenseite (oder das Gericht) müsste all diese Unterlagen manuell sichten, um die Beweise zu prüfen. Schon hier offenbaren sich die inhärenten Risiken: Die KI war nicht in der Lage, alle relevanten Informationen zu filtern \cite{Luther2026}. Während ein geschulter Mensch zwar in der Lage wäre, dies korrekt zu tun, würde er dafür unverhältnismäßig viel Zeit benötigen – ein Nachteil, der in der Praxis oft zu prozessualen Zugeständnissen oder Vergleichsverhandlungen führt.
	
	\subsubsection{Die Notwendigkeit eines \enquote{Sovereign AI}}
	Prof. Steinrötter von der Universität Potsdam weist darauf hin, dass die öffentliche Hand angesichts dieser Übermacht nicht nur die technologischen, sondern auch die rechtlichen Rahmenbedingungen schaffen muss. Er betont: \enquote{Letztlich müssen alle relevanten Gesetze, die dazugehörige Rechtsprechung und die Verwaltungspraxis in das System eingespeist werden.} Der Staat müsse insbesondere auf \enquote{Sovereign AI} – also souveräne, europäische KI-Lösungen – setzen, um nicht von US-Konzernen abhängig zu werden \cite{Steinroetter2026}.
	
	\subsubsection{Reputations- und Haftungsrisiken als Bremsfaktor für Kanzleien}
	Die Kehrseite des Wettbewerbsvorteils ist das immense Risiko bei unsachgemäßer Nutzung. Die aggressive Nutzung von KI kann für Kanzleien nach hinten losgehen, was ihre Zurückhaltung erklären könnte. Das Landgericht Frankfurt am Main rügte im September 2025 einen Anwalt scharf, weil dieser in einem Schriftsatz frei erfundene BGH-Zitate aus einem KI-Tool verwendet hatte (Az. 2-13 S 56/24). Die Richter sprachen von einer \enquote{kompletten Fälschung} \cite{LG Frankfurt2025}. Das Landgericht Darmstadt kürzte in einem Fall das Honorar eines Sachverständigen auf 0,00 Euro, weil er ein maßgeblich von einer KI erstelltes Gutachten nicht deklariert hatte (Az. 19 O 527/16) \cite{Kramarz2026}. In einem anderen Fall wertete das Kammergericht Berlin automatisierte Schriftsätze ohne Bezug zum konkreten Tatvorwurf als dysfunktionale Prozessführung \cite{KG Berlin2025}.
	
	\subsection{Reaktionen der Gerichte und Grenzen des richterlichen Ersatzes}
	
	Der österreichische Oberste Gerichtshof (OGH) wies im Oktober 2025 eine offensichtlich mit KI erstellte Nichtigkeitsbeschwerde zurück, die sich auf erfundene angebliche höchstgerichtliche Entscheidungen stützte. Das Gericht stellte klar, dass dieses Vorbringen dem erforderlichen Argumentationsniveau nicht genüge \cite{OGH2025}.
	
	Der Thüringer Richterbund betonte im August 2025: \enquote{Was aus meiner Sicht gar nicht geht, ist, dass Sie sich von ChatGPT die Urteile schreiben lassen} \cite{Richterbund2025}. Die Findung eines Urteils müsse immer einem Richter vorbehalten bleiben. KI dürfe für Gerichte immer nur ein Hilfsmittel sein, aber \enquote{nie den Richter ersetzen} \cite{Richterbund2025}.
	
	\subsection{Alltägliche Gefahren für nicht IT-affine Nutzer}
	
	Während die vorangegangenen Abschnitte die institutionellen Herausforderungen für Polizei und Justiz beleuchtet haben, ist KI längst zu einem Massenphänomen geworden, das den Alltag von Millionen Menschen prägt. Laut der aktuellen TÜV-ChatGPT-Studie 2025 nutzen bereits 65 Prozent der Bundesbürgerinnen und -bürger generative KI-Tools – ein deutlicher Anstieg gegenüber 37 Prozent im Jahr 2023. Besonders die jüngere Generation (91 Prozent der 16–29-Jährigen) setzt regelmäßig auf Tools wie ChatGPT, Google Gemini oder DeepL \cite{TUEV2025}. Die Kehrseite dieser rasanten Verbreitung ist eine gravierende Sicherheitslücke im Umgang mit KI. Viele Nutzerinnen und Nutzer erkennen KI-generierte Inhalte nicht oder unterschätzen die Gefahren durch Deepfakes, fehlerhafte Ergebnisse und Datenmissbrauch \cite{TUEV2025}. Dieser Abschnitt widmet sich den spezifischen Risiken für genau diese nicht IT-affinen Nutzer.
	
	\subsubsection{Cyberkriminalität: Neue Dimensionen von Betrug und Malware}
	
	KI eröffnet Kriminellen völlig neue Angriffsmöglichkeiten. Eine aktuelle Bitdefender-Studie gibt Einblicke in die zentralen Verhaltensweisen, Praktiken und Bedenken von Verbrauchern weltweit im Bereich der Cybersicherheit \cite{Bitdefender2025}. Besonders alarmierend ist die Entwicklung von KI-gestützten Phishing-Angriffen: KI verbessert Phishing-Mails, Smishing-Nachrichten und Deepfake-Anrufe so massiv, dass sie täuschend echt wirken und klassische Indikatoren für Betrug kaum noch greifen \cite{SecurityInsider2025a}. Drei von vier Befragten geben an, heute mehr Sorge um die Sicherheit ihrer persönlichen Daten zu haben als noch vor fünf Jahren. Besonders besorgniserregend: Nur 23 Prozent fühlen sich in der Lage, legitime von betrügerischen Inhalten sicher zu unterscheiden \cite{SecurityInsider2025b}. Als größte Bedrohung nennen 39 Prozent der Verbraucher KI-gestütztes Phishing \cite{SecurityInsider2025b}.
	
	Neben Phishing hat sich insbesondere das Voice Cloning als gefährliche Betrugsmasche etabliert. Modelle können Stimmen bereits nach wenigen Sekunden nachahmen. Angreifer imitieren vertrauenswürdige Personen, wie Vorgesetzte oder Familienmitglieder, um ihre Opfer zu Überweisungen, zur Weitergabe von Zugangsdaten oder zur Freigabe sensibler Daten zu verleiten \cite{Watchguard2025}. Die Erkennung solcher Manipulationen kommt häufig zu spät, sodass Angriffe in der Regel bereits erfolgreich sind, bevor jemand reagieren kann \cite{Watchguard2025}.
	
	Auch die Malware-Entwicklung hat durch KI eine neue Qualität erreicht. Im Mai 2025 wurde eine polymorphe KI-Malware bekannt, die anders als klassische polymorphe Varianten nicht nur innerhalb vorgegebener Muster variiert, sondern durch generative KI bei jeder Ausführung nahezu einzigartig ist \cite{ConnectProfessional2025}. Angriffsframeworks basieren zunehmend auf KI und passen ihre Payloads in Echtzeit an. Ein Proof-of-Concept namens \enquote{Ransomware 3.0} hat demonstriert, wie ein KI-gesteuerter Ransomware-Prototyp Angriffsphasen dynamisch planen, adaptieren und ausführen kann \cite{Watchguard2025}.
	
	Eine weitere Studie der Elon University aus dem Jahr 2025 zeigt, dass 58 Prozent der KI-Nutzenden regelmäßig mehrere Modelle parallel einsetzen. Laut Daten des Unternehmens Harmonic erfolgen 45,4 Prozent der Eingaben sensibler Informationen über persönliche Konten – ohne Kontrolle durch die IT-Abteilung \cite{Itwelt2025}. Noch besorgniserregender: 44 Prozent der Befragten gaben an, bei der Nutzung von KI gegen Unternehmensrichtlinien verstoßen zu haben. Risiken entstehen zudem durch sogenannte Schatten-KI, wenn Mitarbeitende private Accounts nutzen und so die Kontrolle durch das Unternehmen entfällt \cite{Datenschutz2025}. Die zunehmende Abhängigkeit von KI-Systemen birgt auch potenzielle Risiken durch unzureichende Datenschutz- und Sicherheitsmaßnahmen \cite{Europarl2025}.
	
	\subsubsection{Fehlinformation und Deepfakes: Die Unterwanderung der digitalen Realität}
	
	Gezielte Einflussnahme auf Wahlen sowie inzwischen alltägliche Desinformation und Fake News sind ernste Bedrohungen für unsere Demokratie \cite{Tagesschau2025}. Kriminelle und politische Akteure können mithilfe von KI-Klonen glaubwürdige Deepfakes erstellen und damit betrügerische Absichten verfolgen \cite{Tagesschau2025}. Social Bots und Algorithmen manipulieren unsere Wahrnehmung, während KI-generierte Stimmen oft kaum noch als solche zu erkennen sind \cite{Sonntagsblatt2025}.
	
	Die IU Studie 2025 zeigt, dass Fake News und Deepfakes konkrete Auswirkungen auf das (Sozial-)Leben und die Wahlentscheidungen der Befragten haben \cite{IU2025}. KI-generierte Inhalte werden oft für echt gehalten. Betrüger können generative KI-Tools nutzen, um raffinierte und effektive Phishing-E-Mails zu erstellen, die ihnen helfen, vertrauliche Informationen zu erhalten \cite{IBM2025}. Im März 2025 setzte sich ein besorgniserregendes Muster fort: In 33 Prozent der Fälle wiederholten KI-Modelle Falschbehauptungen, die über das \enquote{Pravda}-Netzwerk verbreitet wurden \cite{NewsGuard2025}.
	
	\subsubsection{Dark Patterns: Wenn KI zur Manipulationsmaschine wird}
	
	Eine aktuelle Studie des Nationalen Forschungszentrums für angewandte Cybersicherheit ATHENE, der University of Glasgow und der Humboldt-Universität zu Berlin zeigt, dass KI-Anwendungen wie ChatGPT, Gemini und Claude AI selbst bei neutralen Anfragen regelmäßig manipulative Elemente wie gefälschte Kundenrezensionen, künstlich erzeugte Zeitdruck-Anzeigen oder irreführende Preisvergleiche in den HTML-Code von Webseiten einbauen – ohne darauf hinzuweisen, dass es sich um rechtlich und ethisch fragwürdige Elemente handelt \cite{ATHENE2025}. Dr. Veronika Krauß, eine der Autorinnen der Studie, erklärt: \enquote{Diese Designmuster sind darauf ausgelegt, Nutzende zu Handlungen zu bewegen, die sie eigentlich nicht beabsichtigen. Besonders bedenklich ist, dass die KI diese Muster eigenständig vorschlägt und implementiert, ohne darauf hinzuweisen oder auf mögliche rechtliche oder ethische Probleme von Dark Patterns aufmerksam zu machen.} \cite{ATHENE2025}
	
	Die Forschenden identifizierten unter anderem vorausgewählte Produkte im Warenkorb, gefälschte euphorische Kundenrezensionen sowie künstlich erzeugte Dringlichkeit durch Zeitdruck-Anzeigen. Besorgniserregend ist, dass die KI-Systeme die psychologischen Wirkmechanismen hinter den manipulativen Designs durchaus verstehen und erklären können \cite{ATHENE2025}. Die Ursache liegt in den Trainingsdaten: Im Internet gibt es genug Webseiten, die manipulative Designs nutzen. Daraus lernt die KI, selbst ähnlichen HTML-Code zu erstellen, um bestimmte Zielvorgaben zu erreichen \cite{ATHENE2025}.
	
	Eine weitere subtile, aber nicht weniger gefährliche Manipulationsform ist die sogenannte Sykophantie – die unkritische Zustimmung von KI-Chatbots. Studien bestätigen, dass sich diese Ja-Sagerei nicht wirklich aus Chatbots herausbekommen lässt, da sie als Produkte eine angenehme Nutzungserfahrung bieten sollen \cite{T3N2025}. Im schlimmsten Fall kann uneingeschränkte Zustimmung und Bekräftigung durch ChatGPT und Co. zu sogenannten \enquote{KI-Psychosen} und Abwärtsspiralen führen \cite{T3N2025}. Der Digital Fairness Act (DFA) soll in Zukunft umfassende Vorgaben gegen solche manipulativen Designelemente einführen \cite{BigDataInsider2025}.
	
	\subsubsection{Datenschutz und rechtliche Grauzonen}
	
	Die Bedenken der Verbraucher gegenüber KI sind enorm. 54 Prozent der deutschen Befragten gaben an, besorgt über den Einfluss von KI auf die Sicherheit ihrer digitalen Identitäten zu sein. Jeder zweite Deutsche würde einem KI-Agenten keine persönlichen Daten anvertrauen, und 41 Prozent erledigen Aufgaben lieber selbst, als sich auf Künstliche Intelligenz zu verlassen \cite{Okta2025}.
	
	Eine aktuelle Analyse der Datenschutzfirma Incogni zeigt, dass KI-Anbieter beim Training ihrer Modelle Schranken wie die robots.txt-Datei ignorieren und Daten aus nicht näher spezifizierten Quellen wie Datenbrokern oder Finanzinstituten beziehen \cite{JuraUniSaarland2025}. Das auf Datenschutz spezialisierte Unternehmen Incogni hat im Rahmen des Privacy-Ranking-Reports 2025 neun KI-Anbieter untersucht und kommt zu dem Schluss, dass bei Gemini und Copilot besonders viele Datenschutzrisiken bestehen \cite{SecurityInsider2025c}. 13 Prozent der KI-Nutzer haben bereits persönliche Daten wie Adressen oder Passwörter in KI-Systeme eingegeben, und 50 Prozent äußern große Sorge vor Datenmissbrauch. In Unternehmen fehlt es an klaren Regeln: 54 Prozent der Beschäftigten berichten von keinen Vorgaben oder Verboten für den KI-Einsatz \cite{TUEV2025}.
	
	Ein besonders komplexes Problemfeld ist der urheberrechtliche Status KI-generierter Inhalte. Bei Deepfakes handelt es sich regelmäßig um vollständig neue, KI-generierte Inhalte, die kein bestehendes Werk verändern. Dies gilt etwa für die meisten KI-generierten Songs, Videos oder Bilder. Weder die eigene Stimme noch das äußere Erscheinungsbild genießen urheberrechtlichen Schutz \cite{JuraUniSaarland2025a}. Eine bloße Verwechslungsgefahr kennt das Urheberrecht nicht. Zusammenfassend lässt sich festhalten, dass das Urheberrecht in erster Linie das Werk schützt – nicht die Person. Der Schutz der eigenen Stimme oder des eigenen Abbilds erfolgt, wenn überhaupt, über das allgemeine Persönlichkeitsrecht oder unlauteren Wettbewerb \cite{JuraUniSaarland2025a}. Während in den USA mit dem ELVIS Act ein Schutz vor kommerziellem Missbrauch KI-simulierter Stimmen besteht, fehlen entsprechende Regelungen bislang im deutschen und europäischen Raum \cite{Oppenhoff2025}. Personen, die in Deepfakes imitiert werden, werden in ihrem allgemeinen Persönlichkeitsrecht verletzt. Da die betroffene Person in die Erstellung des Deepfakes sicher nicht eingewilligt hat, ist das Deepfaken meist eine rechtswidrige Verarbeitung \cite{WBS2025}. Dennoch bleiben die rechtlichen Möglichkeiten für Betroffene derzeit oft unzureichend.
	
	\subsubsection{Schutz durch die KI-Verordnung?}
	
	Die europäische KI-Verordnung (KI-VO) soll eine Vielzahl realer KI-Risiken adressieren, denen Verbraucherinnen und Verbraucher täglich ausgesetzt sind \cite{VZBV2025}. Sie räumt Verbraucherinnen und Verbrauchern jedoch nur wenige Individualrechte ein: das Recht auf Beschwerde und das Recht auf Erläuterung einer Hochrisiko-KI-Entscheidung \cite{VZBV2025}. Der EU AI Act ist zudem nur 32 Prozent der Menschen bekannt \cite{TUEV2025}. Seit dem 2. Februar 2025 gelten Verbote für besonders riskante KI-Anwendungen in der EU \cite{Rechtsanwalt2025}. Die KI-VO geht auf die Ziele des europäischen Verbraucherschutzes jedoch nur an sehr wenigen Stellen ein \cite{CRonline2025}. Ein breites Bündnis warnt vor einer möglichen Aushöhlung der KI-Regeln durch geplante \enquote{Vereinfachungen} \cite{Netzpolitik2026}. Für KI-Systeme, die für Verbraucherinnen und Verbraucher besonders riskant sind, gelten strenge Qualitäts- und Verfahrensanforderungen. Dazu zählen solche, die zur Kreditvergabe oder zur Prämienfestsetzung im Versicherungsbereich eingesetzt werden \cite{BMJV2025}.
	
	\section{Die Rolle der deutschen und europäischen IT-Industrie und der Open-Source-Community}
	
	Angesichts der skizzierten Herausforderungen stellt sich zwangsläufig die Frage, inwieweit die öffentliche Hand die notwendige technologische Souveränität überhaupt erreichen kann, ohne sich noch tiefer in Abhängigkeiten von großen US-Tech-Konzernen zu begeben. Dies erfordert einen differenzierten Blick auf die Strategien der Bundesregierung, der europäischen IT-Wirtschaft und der Open-Source-Community.
	
	\subsection{Die Open-Source-Strategie des Bundes: Das Zentrum für Digitale Souveränität (ZenDiS)}
	
	Um die öffentliche Verwaltung aus der Abhängigkeit von wenigen großen Tech-Konzernen zu befreien, wurde im Dezember 2022 das \textbf{Zentrum für Digitale Souveränität der Öffentlichen Verwaltung (ZenDiS)} gegründet \cite{ZenDiS2025}. Das ZenDiS ist eine GmbH im Besitz des Bundes und verfolgt das Ziel, Bund, Ländern und Kommunen ein umfassendes Souveränitätspaket aus Plattform, Produkten und Beratung zur Verfügung zu stellen \cite{ZenDiS2025}. Kern der Strategie ist das Prinzip \enquote{Public Money, Public Code}, das im Koalitionsvertrag 2021 verankert wurde und besagt, dass mit öffentlichen Geldern entwickelte Software auch der Allgemeinheit als Open Source zur Verfügung stehen soll \cite{SovereignTechFund2025}.
	
	Das ZenDiS hat zwei zentrale Plattformen geschaffen:
	\begin{itemize}
		\item \textbf{openCode:} Eine offene Plattform (basierend auf GitLab) zur Entwicklung, Bereitstellung und Nachnutzung von Open-Source-Software für die öffentliche Verwaltung. Sie dient als zentrale Anlaufstelle für die kollaborative Softwareentwicklung von Bund, Ländern und Kommunen. Mittlerweile basieren über 4.000 Projekte auf openCode \cite{Heise2025a}.
		\item \textbf{openDesk:} Eine souveräne Open-Source-Arbeitsplatzlösung, die als vollwertige Alternative zu Microsoft 365 und Teams konzipiert ist. openDesk integriert etablierte Open-Source-Tools wie Open-Xchange (Groupware), OpenProject (Projektmanagement), Nextcloud (Dateisynchronisation), XWiki (Wissensmanagement), Matrix und Jitsi (Kommunikation) sowie Collabora (Web-Office) in einer einheitlichen Plattform \cite{AllThingsOpen2025}. Die Verfügbarkeit als Source-Code, Helm-Charts und bald auch als SaaS-Angebot macht es zu einem echten Werkzeug für die digitale Souveränität \cite{AllThingsOpen2025}.
	\end{itemize}
	
	\subsection{Die Open-Source-Community als Fundament: Der Sovereign Tech Fund}
	
	Neben der Bereitstellung fertiger Lösungen hat die Bundesregierung ein international vielbeachtetes Instrument zur gezielten Stärkung der Open-Source-Infrastruktur geschaffen: den \textbf{Sovereign Tech Fund (STF)}. Der STF ist ein Förderprogramm des Bundesministeriums für Wirtschaft und Klimaschutz (BMWK), das darauf abzielt, die Sicherheit, Wartung und Resilienz kritischer Open-Source-Basistechnologien zu stärken \cite{SovereignTechFund2025}. Seit seiner Gründung im Jahr 2022 hat der STF über 24,6 Millionen Euro bereitgestellt, um mehr als 60 Open-Source-Projekte weltweit zu unterstützen \cite{SovereignTechFund2025}. Die Mindestförderung pro Projekt beträgt 50.000 Euro, in der Praxis wurden bereits Investitionen von bis zu einer Million Euro getätigt \cite{SovereignTechFund2025}.
	
	Der STF adressiert damit die strukturelle Herausforderung, dass viele grundlegende Open-Source-Komponenten – wie Bibliotheken, Protokolle und Entwicklungswerkzeuge – aufgrund ihrer Allgegenwart in der digitalen Infrastruktur oft vernachlässigt werden, obwohl sie das Rückgrat der digitalen Wirtschaft und Verwaltung bilden. Mit Programmen wie dem \enquote{Sovereign Tech Fellowship}, das Maintainer*innen von wichtigen Open-Source-Komponenten für ihre Arbeit bezahlt, verfolgt der STF einen ganzheitlichen Ansatz zur Sicherung des Open-Source-Ökosystems im öffentlichen Interesse \cite{STFNews2025}.
	
	\subsection{Deutsche IT-Wirtschaft: Zwischen Eigenentwicklung und Abhängigkeit}
	
	Die deutsche IT-Wirtschaft ist in diesem Ökosystem auf verschiedene Weisen aktiv. Einerseits gibt es vielversprechende Eigenentwicklungen, die bewusst auf Open Source setzen, um Unabhängigkeit zu gewährleisten.
	
	Ein herausragendes Beispiel ist der KI-Assistent \textbf{F13}, der vom Innovationslabor Baden-Württemberg entwickelt und im Juli 2025 als Open-Source-Software (MPL 2.0 Lizenz) auf openCode veröffentlicht wurde \cite{Staatsanzeiger2025}. F13 ist eine modellagnostische, datensouveräne KI-Assistenz, die auf deutschen Servern betrieben wird und somit DSGVO-konform ist. Sie umfasst bereits Funktionen wie Chat, Dokumentenzusammenfassung und Recherche in verwaltungsrelevanten Quellen und wird bereits von der Polizei Baden-Württemberg genutzt \cite{Staatsanzeiger2025}. Das ZenDiS hat das Projekt bei der Lizenzprüfung und der Community-Strategie beratend unterstützt \cite{ZenDiS2025b}. Diese Entwicklung zeigt, dass eine datenschutzkonforme und souveräne KI-Infrastruktur für die öffentliche Hand technisch möglich ist – wenn der politische Wille und die entsprechenden Ressourcen vorhanden sind.
	
	Andererseits wird die anhaltende Abhängigkeit von proprietären US-Lösungen besonders deutlich im Bereich der Polizei-Datenanalyse. Die \textbf{Palantir}-Software ist in den Polizeibehörden von Bayern, Hessen, Baden-Württemberg und Nordrhein-Westfalen im Einsatz \cite{Tagesschau2025}. Politisch ist man sich einig, dass eine souveräne Alternative geschaffen werden muss. So bekräftigte die Innenministerkonferenz im Juni 2025 das Ziel, eine \enquote{Verfahrensübergreifende Recherche- und Analyseplattform} zu schaffen, die frei von außereuropäischen Einflüssen ist \cite{Tagesschau2025}.
	
	Dennoch scheiterten vielversprechende Eigeninitiativen. So schloss sich 2023 ein Konsortium aus sieben deutschen Unternehmen (Start-Ups, Mittelständler und ein Großkonzern) unter dem Namen \textbf{NaSa} zusammen, um eine konkrete Palantir-Alternative zu entwickeln. Sachar Paulus, Professor für IT-Sicherheit an der Hochschule Mannheim, hielt eine solche Lösung für realistisch und schätzte, dass innerhalb von zwölf bis 18 Monaten ein Prototyp umsetzbar sei \cite{Tagesschau2025}. Trotz eines Zeitplans und der Vorstellung im Innenausschuss des Bundestags fand das Projekt keinen Auftraggeber – weder auf Landes- noch auf Bundesebene, sodass es scheiterte \cite{Tagesschau2025}. Dies zeigt ein strukturelles Problem auf: Während die Politik digitale Souveränität propagiert, scheitert die konkrete Umsetzung oft an fehlenden verbindlichen Aufträgen und der mangelnden Bereitschaft, langfristig in heimische Entwicklungen zu investieren.
	
	\subsection{Europäische Initiativen: Open Source in der Justiz}
	
	Auch im Justizbereich gibt es ermutigende Open-Source-Projekte, die den Weg für mehr Souveränität ebnen. So gewann das Bundesministerium der Justiz und für Verbraucherschutz mit dem Projekt \textbf{„Zugang zum Recht“} den Open-Source-Wettbewerb der Open Source Business Alliance (OSBA) in der Kategorie „Fachverfahren“ \cite{Heise2025c}. Die Anwendung bietet digitale Justizdienste an, allen voran die digitale Klage für Fluggastrechte, die es Bürgern ermöglicht, direkt eine digitale Klage ohne Anwalt einzureichen. Jährlich ließen sich so rund 790.000 Euro an Beratungshilfekosten einsparen \cite{Heise2025c}.
	
	Ein weiteres Leuchtturmprojekt ist \textbf{NeuRIS}, das neue Rechtsinformationssystem des Bundes. Der DigitalService entwickelt gemeinsam mit dem BMJV und dem Bundesamt für Justiz ein Portal, das Gesetze, Gerichtsentscheidungen und Verwaltungsvorschriften des Bundes zentral, barrierefrei und über eine offene API kostenlos zur Verfügung stellt \cite{NeuRIS2025}. Dies ermöglicht nicht nur einen modernen Zugang zum Recht, sondern schafft auch eine datensouveräne Grundlage für Legal-Tech-Anwendungen.
	
	\subsection{Ausblick: Open Source als Standard}
	
	Die Weichen für eine stärkere Open-Source-Nutzung sind gestellt. Die überarbeiteten EVB-IT-Verträge, die im Februar 2026 veröffentlicht wurden, machen Open Source zum Standard für die öffentliche Verwaltung und standardisieren die rechtskonforme Beschaffung von Open-Source-Software \cite{ZenDiS2025}. Auch im EU-weiten Kontext wird das Thema vorangetrieben: Im Oktober 2025 wurde das \enquote{EU Digital Commons Consortium} (DC-EDIC) offiziell gegründet, um digitale Gemeingüte in Europa zu fördern. Deutschland wird hier durch das Bundesministerium für Digitales und Staatsmodernisierung vertreten \cite{ZenDiS2025}.
	
	\section{Verbindende Querschnittsthemen}
	
	\subsection{Datenschutz und Grundrechte}
	
	Die automatisierte Analyse personenbezogener Daten durch KI stellt einen massiven Eingriff in die informationelle Selbstbestimmung dar. Die DSK betont, dass automatisierte Datenanalysen grundsätzlich alle Menschen betreffen können, ohne dass sie durch ihr Verhalten einen Anlass für polizeiliche Ermittlungen gegeben hätten \cite{DSK2025}.
	
	Besonders problematisch sind laut DSK Datenanalysen von Personen, die keinen Anlass für Ermittlungen gegeben haben, sowie Zweckänderungen \cite{Datenschutzticker2025}. Die Humanistische Union warnt, dass KI-gestützte \enquote{Superdatenbanken} im Widerspruch zur EU-KI-Verordnung stehen und weitreichende Profilbildung auch von Unbeteiligten ermöglichen \cite{HumanUnion2025}.
	
	\subsection{EU-KI-VO als neuer Rechtsrahmen}
	
	Die EU-KI-Verordnung (KI-VO) reguliert KI-Systeme nach einem risikobasierten Ansatz und unterwirft Hochrisiko-KI-Systeme besonders strengen Anforderungen \cite{Kriminalpolizei2026}. Biometrische Massenüberwachung in Echtzeit ist nach Art. 5 KI-VO grundsätzlich eine verbotene Praxis \cite{RAV2025}.
	
	Die Bundesregierung plant jedoch offenbar eine Auslagerung des biometrischen Abgleichs ins nicht-europäische Ausland, um diese Verbote zu umgehen \cite{Gruene2026}. Innenminister Dobrindt fordert sogar eine \enquote{Korrektur} der Leitlinien für Hochrisiko-Systeme \cite{Dobrindt2026}. Die KI-VO erstreckt sich nach überwiegender Auffassung auch auf präventiv polizeiliche Aufgabenwahrnehmung \cite{Kriminalpolizei2026}.
	
	\section{Fazit}
	
	KI für Polizei und Justiz stellt kein eindimensionales Problem dar, sondern eröffnet ein breites Spektrum neuer Herausforderungen. Die größten Hürden liegen im Spannungsfeld zwischen technologischen Möglichkeiten und rechtsstaatlichen Prinzipien: bei der Frage nach angemessenen rechtlichen Rahmenbedingungen, dem Schutz vor Diskriminierung, der Wahrung richterlicher Unabhängigkeit und der Sicherung informationeller Selbstbestimmung. Besonders alarmierend sind die konkreten Gefahren von Fehlurteilen durch KI-Halluzinationen, manipulierte Beweise und verzerrte Entscheidungsgrundlagen, die bereits jetzt zu nachweisbaren Verfahrensfehlern führen. Hinzu tritt ein neuartiges strukturelles Problem: Private Kanzleien nutzen KI strategisch zu ihrem Vorteil, was zu einem gefährlichen Ungleichgewicht im Rechtsstaat führen kann, wenn der Staat technologisch nicht Schritt hält.
	
	Die strategische Antwort des Bundes auf diese Herausforderung – ZenDiS, STF, openCode, openDesk, F13 und NeuRIS – zeigt, dass die technologischen und konzeptionellen Grundlagen für eine souveräne digitale Verwaltung vorhanden sind. Die entscheidende Frage ist nun, ob der politische Wille ausreicht, um diese Werkzeuge konsequent einzusetzen und die Entwicklung durch verbindliche Ausschreibungen und langfristige Investitionen zu fördern. Nur so kann Deutschland die drohende digitale Spaltung in eine \enquote{Zwei-Klassen-Justiz} verhindern und die Fairness des Rechtsstaats auch im Zeitalter der KI bewahren.
	
	\newpage
	\begin{thebibliography}{99}
		
		% ... (alle bisherigen Quellen wie zuvor)
		
		% Neue Quellen für Abschnitt 3.8
		\bibitem{TUEV2025} TÜV-Verband (2025): \emph{ChatGPT-Studie 2025: Wie Deutschland generative KI nutzt}. November 2025. S. 3-42.
		
		\bibitem{Bitdefender2025} Bitdefender (2025): \emph{2025 Consumer Cybersecurity Survey}. November 2025.
		
		\bibitem{SecurityInsider2025a} Security-Insider (2025): \emph{Mobile Security Index 2025: Risiken durch KI und Anwenderfehler}. 29. Dezember 2025.
		
		\bibitem{SecurityInsider2025b} Security-Insider (2025): \emph{Vertrauensverlust in Künstlicher Intelligenz: Sicherheitsbedenken und Ruf nach Regulierung}. 6. November 2025.
		
		\bibitem{Watchguard2025} WatchGuard Technologies (2025): \emph{Die fünf größten Sicherheitsrisiken durch KI}. 7. Oktober 2025.
		
		\bibitem{ConnectProfessional2025} Connect Professional (2025): \emph{Das Dilemma generativer KI}. 17. November 2025.
		
		\bibitem{Itwelt2025} Itwelt.at (2025): \emph{Sicherheitsrisiko KI: Großteil der Tools bereits von Datenpannen betroffen}. 12. Mai 2025.
		
		\bibitem{Datenschutz2025} Dr. Datenschutz (2025): \emph{Shortnews im Juni 2025 – KW24}. 30. November 2025.
		
		\bibitem{Europarl2025} Europäisches Parlament (2025): \emph{Künstliche Intelligenz: Chancen und Risiken}. 2. Mai 2025.
		
		\bibitem{Tagesschau2025a} Tagesschau (2025): \emph{Mithilfe von KI gefälschte Videos: Das Ende der Wirklichkeit?} 6. Juni 2025.
		
		\bibitem{Sonntagsblatt2025} Sonntagsblatt (2025): \emph{Social Bots und Algorithmen: Wie KI die Verbreitung von Fake News revolutioniert}. 23. Juni 2025.
		
		\bibitem{IU2025} IU Internationale Hochschule (2025): \emph{Gefahren von Fake News und Deepfakes – IU Studie}. 4. September 2025.
		
		\bibitem{IBM2025} IBM (2025): \emph{KI-Fehlinformationen}. 2025.
		
		\bibitem{NewsGuard2025} NewsGuard (2025): \emph{Anteil falscher Antworten von KI-Chatbots hat sich in einem ...} März 2025.
		
		\bibitem{ATHENE2025} ATHENE (2025): \emph{Studie zeigt: KI erzeugt manipulative Designmuster in Webseiten}. 30. April 2025.
		
		\bibitem{T3N2025} t3n (2026): \emph{KI-Chatbots und Sykophantie: So stark ist die Bauchpinselei bei Claude, Gemini und ChatGPT}. 18. April 2026.
		
		\bibitem{BigDataInsider2025} BigData-Insider (2026): \emph{Digital Fairness Act (DFA): Was Unternehmen jetzt vorbereiten sollten}. 26. Mai 2026.
		
		\bibitem{Okta2025} Okta (2025): \emph{Customer Identity Trends Report 2025}. 11. Juli 2025.
		
		\bibitem{JuraUniSaarland2025} Jura Uni Saarland (2025): \emph{Datenschutz unter Druck: Wie KI-Anbieter mit persönlichen Daten umgehen}. 30. November 2025.
		
		\bibitem{SecurityInsider2025c} Security-Insider (2025): \emph{ChatGPT, Gemini und Co. unter der Lupe: Welche KI schnüffelt am meisten?} 2. September 2025.
		
		\bibitem{JuraUniSaarland2025a} Jura Uni Saarland (2025): \emph{KI und Urheberrecht: Deepfakes}. November 2025.
		
		\bibitem{Oppenhoff2025} Oppenhoff (2025): \emph{Der Schutz der eigenen Stimme in Zeiten von KI Voice Cloning}. 3. Oktober 2025.
		
		\bibitem{WBS2025} WBS Legal (2025): \emph{KI, "Deepfakes" und Face-Swap-Apps können Personen imitieren}. 20. November 2025.
		
		\bibitem{VZBV2025} Verbraucherzentrale Bundesverband (vzbv) (2025): \emph{Stellungnahme zum Entwurf des KI-VO-Durchführungsgesetzes}. 8. Oktober 2025.
		
		\bibitem{Rechtsanwalt2025} rechtsanwalt.com (2026): \emph{KI-Haftung: EU-KI-Gesetz trifft deutsche Verbraucher}. 15. Mai 2026.
		
		\bibitem{CRonline2025} CRonline (2024): \emph{Die Durchsetzung der KI-Verordnung im Rahmen des digitalen Schuldrechts}. 2024.
		
		\bibitem{Netzpolitik2026} netzpolitik.org (2026): \emph{Breites Bündnis warnt vor verwässerten KI-Regeln}. 8. April 2026.
		
		\bibitem{BMJV2025} Bundesministerium der Justiz und für Verbraucherschutz (2025): \emph{Verbraucherschutz beim Einsatz von Künstlicher Intelligenz}. 2025.
		
		% ... (alle weiteren Quellen wie zuvor)
		
	\end{thebibliography}
	
\end{document}